
\medskip

On considère le cube ABCDEFGH.

On note I le milieu du segment[EH] et on considère le triangle CFI.

L'espace est muni du repère orthonormé $\left(\text{A}~;~ \vect{\text{AB}},\: \vect{\text{AD}},\: \vect{\text{AE}}\right)$ et on admet que le point I a
pour coordonnées $\left(0~;~\dfrac12;~1\right)$ dans ce repère.

\begin{center}
\psset{unit=1cm}
\begin{pspicture}(7.4,8)
\pspolygon(0.2,0.8)(4.6,0.2)(4.6,4.6)(0.2,5.2)%BCGF
\psline(4.6,0.2)(7.2,2.6)(7.2,7)(4.6,4.6)%CDHG
\psline(7.2,7)(2.8,7.6)(0.2,5.2)%HEF
\psline[linestyle=dashed](0.2,0.8)(2.8,3.2)(7.2,2.6)%BAD
\psline[linestyle=dashed](2.8,3.2)(2.8,7.6)%AE
\psline(5,7.3)(0.2,5.2)(4.6,0.2)%IFC
\psline[linestyle=dashed](5,7.3)(4.6,0.2)%IC
\psdots(2.8,3.2)(0.2,0.8)(4.6,0.2)(7.2,2.6)(2.8,7.6)(0.2,5.2)(4.6,4.6)(7.2,7)(5,7.3)(2.4,4.9)
\uput[l](2.8,3.2){A} \uput[l](0.2,0.8){B} \uput[dr](4.6,0.2){C} \uput[r](7.2,2.6){D}
\uput[u](2.8,7.6){E} \uput[l](0.2,5.2){F} \uput[dl](4.6,4.6){G} \uput[ur](7.2,7){H}
\uput[u](5,7.3){I} \uput[u](2.4,4.9){J}
\end{pspicture}
\end{center}
\smallskip

\begin{enumerate}
\item 
	\begin{enumerate}
		\item Donner sans justifier les coordonnées des points C, F et G.
		\item Démontrer que le vecteur $\vect{n}\begin{pmatrix}1\\2\\2\end{pmatrix}$ est normal au plan (CFI).
		\item Vérifier qu'une équation cartésienne du plan (CFI) est : $x + 2y + 2z - 3 = 0$.
	\end{enumerate}	
\item  On note $d$ la droite passant par G et orthogonale au plan (CFI).
	\begin{enumerate}
		\item Déterminer une représentation paramétrique de la droite $d$.
		\item Démontrer que le point K$\left(\dfrac79~;~\dfrac59~;~\dfrac59\right)$ est le projeté orthogonal du point G sur le plan (CFI).
		\item Déduire des questions précédentes que la distance du point G au plan (CFI) est
égale à $\dfrac23$.
	\end{enumerate}
\item On considère la pyramide GCFI.

\emph{On rappelle que le volume $V$ d'une pyramide est donné par la formule} 

\[V = \frac13 \times b \times h,\]

\emph{où $b$ est l'aire d'une base et $h$ la hauteur associée à cette base}.
	\begin{enumerate}
		\item Démontrer que le volume de la pyramide GCFI est égal à $\dfrac16$, exprimé en unité de volume.
		\item En déduire l'aire du triangle CFI, en unité d'aire.
	\end{enumerate}
\end{enumerate}

%#1 lettre désignant la suite, par défaut u 
%#2 indice initial, par défaut 0 
%#3 valeur du terme initial 
%#4 expression de f 
%#5 nombre de termes 
\newcounter{INDICE} 
\newcommandtwoopt{\SuiteRec}[5][u][0]{% 
\getcoor % récupère les coordonnées de la fenêtres, merci JCC 
\infixtoRPN{#4} 
\psplot{\xmin}{\xmax}{\RPN} 
\SpecialCoor 
%\psaxes{->}(0,0)(\xmin,\ymin)(\xmax,\ymax) 
%\psplot[linewidth=1pt,algebraic,VarStep,VarStepEpsilon=0.0001]{\xmin}{\xmax}{#4} 
\psline[linestyle=dashed](\xmin,\xmin)(\xmax,\xmax)  
\pstVerb{ /x #3 def } 
\psline[ArrowInside=->,ArrowInsideNo=1](! x 0 )(! x \RPN ) 
\psline[ArrowInside=->,ArrowInsideNo=1](! x \RPN )(! \RPN\space dup) 
\psline{-}(! x 0.03 )(! x 0.03 neg) 
\rput(!x 0.08 neg ){$#1_{#2}$} 
\pstVerb{ /x \RPN\space def} 
\setcounter{INDICE}{#2} 
\addtocounter{INDICE}{1} 
\multido{\I=\theINDICE+1}{#5}{% 
\psline{-}(! x 0.03 )(! x 0.03 neg) 
\rput(!x 0.08 neg ){$#1_{\I}$} 
\psline[linestyle=dotted](! x 0 )(! x x) 
\psline[ArrowInside=->,ArrowInsideNo=1](! x x )(! x \RPN ) 
\psline[ArrowInside=->,ArrowInsideNo=1](! x \RPN )(! \RPN\space dup) 
\pstVerb{ /x \RPN\space def} 
}
}
\newcommandtwoopt{\SuiteRecbis}[5][u][0]{% 
\getcoor % récupère les coordonnées de la fenêtres, merci JCC 
\infixtoRPN{1.9 x mul 1 x sub mul} 
\SpecialCoor 
\psline[linestyle=dashed](\xmin,\xmin)(\xmax,\xmax)  
\pstVerb{ /x #3 def } 
\psline[ArrowInside=->,ArrowInsideNo=1](! x 0 )(! x \RPN ) 
\psline[ArrowInside=->,ArrowInsideNo=1](! x \RPN )(! \RPN\space dup) 
\psline{-}(! x 0.03 )(! x 0.03 neg) 
\rput(!x 0.08 neg ){$#1_{#2}$} 
\pstVerb{ /x \RPN\space def} 
\setcounter{INDICE}{#2} 
\addtocounter{INDICE}{1} 
\multido{\I=\theINDICE+1}{4}{% 
\psline{-}(! x 0.03 )(! x 0.03 neg) 
\rput(!x 0.08 neg ){$#1_{\I}$} 
\psline[linestyle=dotted](! x 0 )(! x x) 
\psline[ArrowInside=->,ArrowInsideNo=1](! x x )(! x \RPN ) 
\psline[ArrowInside=->,ArrowInsideNo=1](! x \RPN )(! \RPN\space dup) 
\pstVerb{ /x \RPN\space def} 
}
}